\FigureLayoutDeclare{img/2019/zhejiang/q19_fig1.png}{11.0cm}{40bp 22bp 65bp 43bp}
\FigureLayoutDeclare{img/2019/zhejiang/q21_fig1.png}{6.5cm}{40bp 33bp 7bp 11bp}
\FigureLayoutDeclare{img/2019/zhejiang/q8_fig1.png}{6.5cm}{0bp 0bp 0bp 0bp}

\chapter{2019年浙江卷}

\section{单选题}

\begin{problem}
已知全集 \(U=\{-1,0,1,2,3\}\)，集合 \(A=\{0,1,2\}\)，\(B=\{-1,0,1\}\)，则 \(\left(\complement_U A\right)\cap B=\)（\quad）

\choices
    {\(\{-1\}\)}
    {\(\{0,1\}\)}
    {\(\{-1,2,3\}\)}
    {\(\{-1,0,1,3\}\)}
\end{problem}

\begin{problem}
渐近线方程为 \(x\pm y=0\) 的双曲线的离心率是（\quad）

\choices
    {\(\dfrac{\sqrt2}{2}\)}
    {\(1\)}
    {\(\sqrt2\)}
    {\(2\)}
\end{problem}

\begin{problem}
若实数 \(x\)，\(y\) 满足约束条件 \(\begin{cases}
x-3y+4\ge0,\\
3x-y-4\le0,\\
x+y\ge0,
\end{cases}\) 则 \(z=3x+2y\) 的最大值是（\quad）

\choices
    {\(-1\)}
    {\(1\)}
    {\(10\)}
    {\(12\)}
\end{problem}

\begin{problem}
祖暅是我国南北朝时代的伟大科学家. 他提出的“幂势既同，则积不容易”称为祖暅原理，利用该原理可以得到柱体体积公式 \(V=Sh\)，其中 \(S\) 是柱体的底面积，\(h\) 是柱体的高，若某柱体的三视图如图所示（单位：\(\mathrm{cm}\)），则该柱体的体积（单位：\(\mathrm{cm}^3\)）是（\quad）

\bitmapfigure[max width=0.38\textwidth]{img/2019/zhejiang/q4_fig1.png}

\choices
    {\(158\)}
    {\(162\)}
    {\(182\)}
    {\(324\)}
\end{problem}

\begin{problem}
若 \(a>0\)，\(b>0\)，则“\(a+b\le4\)”是“\(ab\le4\)”的（\quad）

\choices
    {充分不必要条件}
    {必要不充分条件}
    {充分必要条件}
    {既不充分也不必要条件}
\end{problem}

\begin{problem}
在同一直角坐标系中，函数 \(y=\dfrac{1}{a^x}\)，\(y=\log_a\!\left(x+\dfrac12\right)\)（\(a>0\) 且 \(a\neq1\)）图象可能是（\quad）

\choices
    {\(\choicebitmap[width=0.15\paperwidth]{img/2019/zhejiang/q6_choice_a.png}\)}
    {\(\choicebitmap[width=0.15\paperwidth]{img/2019/zhejiang/q6_choice_b.png}\)}
    {\(\choicebitmap[width=0.15\paperwidth]{img/2019/zhejiang/q6_choice_c.png}\)}
    {\(\choicebitmap[width=0.15\paperwidth]{img/2019/zhejiang/q6_choice_d.png}\)}
\end{problem}

\begin{problem}
设 \(0<a<1\)，则随机变量 \(X\) 的分布列如图：

\bitmapfigure[max width=0.32\textwidth]{img/2019/zhejiang/q7_fig1.png}

则当 \(a\) 在 \((0,1)\) 内增大时，\(D(X)\)（\quad）

\choices
    {增大}
    {减小}
    {先增大后减小}
    {先减小后增大}
\end{problem}

\begin{problem}
设三棱锥 \(V-ABC\) 的底面是正三角形，侧棱长均相等，\(P\) 是棱 \(VA\) 上的点（不含端点），记直线 \(PB\) 与直线 \(AC\) 所成角为 \(\alpha\)，直线 \(PB\) 与平面 \(ABC\) 所成角为 \(\beta\)，二面角 \(P-AC-B\) 的平面角为 \(\gamma\)，则（\quad）

\bitmapfigure[max width=0.36\textwidth]{img/2019/zhejiang/q8_fig1.png}

\choices
    {\(\beta<\gamma\)，\(\alpha<\gamma\)}
    {\(\beta<\alpha\)，\(\beta<\gamma\)}
    {\(\beta<\alpha\)，\(\gamma<\alpha\)}
    {\(\alpha<\beta\)，\(\gamma<\beta\)}
\end{problem}

\begin{problem}
已知 \(a\)，\(b\in\R\)，函数 \(f(x)=\begin{cases}
x, & x<0,\\
\dfrac13x^3-\dfrac12(a+1)x^2+ax, & x\ge0,
\end{cases}\) 若函数 \(y=f(x)-ax-b\) 恰有三个零点，则（\quad）

\choices
    {\(a<-1\)，\(b<0\)}
    {\(a<-1\)，\(b>0\)}
    {\(a>-1\)，\(b>0\)}
    {\(a>-1\)，\(b<0\)}
\end{problem}

\begin{problem}
设 \(a\)，\(b\in\R\)，数列 \(\{a_n\}\) 中，\(a_1=a\)，\(a_{n+1}=a_n^2+b\)（\(n\in\N^*\)），则（\quad）

\choices
    {当 \(b=\dfrac12\) 时，\(a_{10}>10\)}
    {当 \(b=\dfrac14\) 时，\(a_{10}>10\)}
    {当 \(b=-2\) 时，\(a_{10}>10\)}
    {当 \(b=-4\) 时，\(a_{10}>10\)}
\end{problem}

\section{填空题}

\begin{problem}
复数 \(z=\dfrac{1}{1+\i}\)（\(\i\) 为虚数单位），则 \(|z|=\fillinblank{}\).
\end{problem}

\begin{problem}
已知圆 \(C\) 的圆心坐标是 \((0,m)\)，半径长是 \(r\). 若直线 \(2x-y+3=0\) 与圆相切于点 \(A(-2,-1)\)，则 \(m=\fillinblank{}\)，\(r=\fillinblank{}\).
\end{problem}

\begin{problem}
在二项式 \((\sqrt2+x)^9\) 的展开式中，常数项是\fillinblank{}；系数为有理数的项的个数是\fillinblank{}.
\end{problem}

\begin{problem}
在 \(\triangle ABC\) 中，\(\angle ABC=90^\circ\)，\(AB=4\)，\(BC=3\)，点 \(D\) 在线段 \(AC\) 上，若 \(\angle BDC=45^\circ\)，则 \(BD=\fillinblank{}\)；\(\cos\angle ABD=\fillinblank{}\).
\end{problem}

\begin{problem}
已知椭圆 \(\dfrac{x^2}{9}+\dfrac{y^2}{5}=1\) 的左焦点为 \(F\)，点 \(P\) 在椭圆上且在 \(x\) 轴的上方，若线段 \(PF\) 的中点在以原点 \(O\) 为圆心，\(|OF|\) 为半径的圆上，则直线 \(PF\) 的斜率是\fillinblank{}.
\end{problem}

\begin{problem}
已知 \(a\in\R\)，函数 \(f(x)=ax^3-x\)，若存在 \(t\in\R\)，使得 \(\left|f(t+2)-f(t)\right|\le\dfrac23\)，则实数 \(a\) 的最大值是\fillinblank{}.
\end{problem}

\begin{problem}
已知正方形 \(ABCD\) 的边长为 \(1\). 当每个 \(\lambda_i\)（\(i=1\)，\(2\)，\(3\)，\(4\)，\(5\)，\(6\)）都取 \(\pm1\) 时，\(\left|\lambda_1\overrightarrow{AB}+\allowbreak\lambda_2\overrightarrow{BC}+\allowbreak\lambda_3\overrightarrow{CD}+\allowbreak\lambda_4\overrightarrow{DA}+\allowbreak\lambda_5\overrightarrow{AC}+\allowbreak\lambda_6\overrightarrow{BD}\right|\) 的最小值是\fillinblank{}；最大值是\fillinblank{}.
\end{problem}

\section{解答题}

\begin{problem}
设函数 \(f(x)=\sin x\)，\(x\in\R\).

\begin{enumerate}
  \item 已知 \(\theta\in[0,2\pi)\)，函数 \(f(x+\theta)\) 是偶函数，求 \(\theta\) 的值；
  \item 求函数 \(y=\left[f\left(x+\dfrac{\pi}{12}\right)\right]^2+\left[f\left(x+\dfrac{\pi}{4}\right)\right]^2\) 的值域.
\end{enumerate}
\end{problem}

\begin{problem}
如图，已知三棱柱 \(ABC-A_1B_1C_1\)，平面 \(A_1AC_1C\perp\) 平面 \(ABC\)，\(\angle ABC=90^\circ\)，\(\angle BAC=30^\circ\)，\(A_1A=A_1C=AC\)，\(E\)，\(F\) 分别是 \(AC\)，\(A_1B_1\) 的中点.

\begin{enumerate}
  \item 证明：\(EF\perp BC\)；
  \item 求直线 \(EF\) 与平面 \(A_1BC\) 所成角的余弦值.
\end{enumerate}

\bitmapfigure[max width=0.28\textwidth]{img/2019/zhejiang/q19_fig1.png}
\end{problem}

\begin{problem}
设等差数列 \(\{a_n\}\) 的前 \(n\) 项和为 \(S_n\)，\(a_3=4\)，\(a_4=S_3\)，数列 \(\{b_n\}\) 满足：对每个 \(n\in\N^*\)，\(S_n+b_n\)，\(S_{n+1}+b_n\)，\(S_{n+2}+b_n\) 成等比数列.

\begin{enumerate}
  \item 求数列 \(\{a_n\}\)，\(\{b_n\}\) 的通项公式；
  \item 记 \(C_n=\sqrt{\dfrac{a_n}{2b_n}}\)（\(n\in\N^*\)），证明：\(C_1+C_2+\cdots+C_n<2\sqrt n\)，\(n\in\N^*\).
\end{enumerate}
\end{problem}

\begin{problem}
如图，已知点 \(F(1,0)\) 为抛物线 \(y^2=2px\)（\(p>0\)）的焦点，过点 \(F\) 的直线交抛物线于 \(A\)，\(B\) 两点，点 \(C\) 在抛物线上，使得 \(\triangle ABC\) 的重心 \(G\) 在 \(x\) 轴上，直线 \(AC\) 交 \(x\) 轴于点 \(Q\)，且 \(Q\) 在点 \(F\) 右侧. 记 \(\triangle AFG\)，\(\triangle CQG\) 的面积分别为 \(S_1\)，\(S_2\).

\begin{enumerate}
  \item 求 \(p\) 的值及抛物线的标准方程；
  \item 求 \(\dfrac{S_1}{S_2}\) 的最小值及此时点 \(G\) 的坐标.
\end{enumerate}

\bitmapfigure[max width=0.32\textwidth]{img/2019/zhejiang/q21_fig1.png}
\end{problem}

\begin{problem}
已知实数 \(a\neq0\)，设函数 \(f(x)=a\ln x+\sqrt{x+1}\)，\(x>0\).

\begin{enumerate}
  \item 当 \(a=-\dfrac34\) 时，求函数 \(f(x)\) 的单调区间；
  \item 对任意 \(x\in\left[\dfrac{1}{\e^2},+\infty\right)\) 均有 \(f(x)\le\dfrac{\sqrt x}{2a}\)，求 \(a\) 的取值范围.
\end{enumerate}

注：\(\e=2.71828\ldots\) 为自然对数的底数.
\end{problem}
